\section{Discussion}

The rate characterization in \cref{obj:thm:uniform-frontier} gives measurement precision a concrete interpretation for estimating the population causal mean at a prespecified dose. Under the structural model in \cref{obj:def:model-class}, with \(c_{\mathrm{lo}}\le(\kappa+1)2^\kappa\le c_{\mathrm{hi}}\) and fixed smoothness and design-decay exponents in their stated ranges, it describes how sample size and the known Gaussian error scale jointly determine attainable precision. Its application therefore rests on substantive justification for the compact latent-dose support, the weak-design envelope, and the causal identification conditions. Within this model, improving measurement precision can move the experiment between three regimes, with comparable rates throughout the transition neighborhoods specified in \cref{obj:thm:uniform-frontier}.

The design-decay exponent \(\kappa\) measures the scarcity of latent doses near the evaluation point. The envelope in \cref{obj:ass:weak-design} implies that the probability of a sufficiently small dose neighborhood of radius \(h\), within either stratum, has order \(h^{\kappa+1}\). Larger values of \(\kappa\) thus represent weaker local support. Under \(c_{\mathrm{lo}}\le(\kappa+1)2^\kappa\le c_{\mathrm{hi}}\), in the direct regime, this scarcity combines with mean smoothness through the effective exponent \(d=2\beta+\kappa+1\), yielding absolute-error and honest-length orders \(n^{-\beta/d}\). For fixed smoothness, increasing design decay slows this order and raises the measurement-error threshold \(n^{-1/d}\) below which direct estimation is attainable. The wider direct-regime range accompanies a coarser attainable localization scale.

In the intermediate Gaussian regime, measurement precision enters more directly. Under \(c_{\mathrm{lo}}\le(\kappa+1)2^\kappa\le c_{\mathrm{hi}}\), the common uncertainty order is
\[
\left\{\frac{\sigma}{\sqrt{\log(e+n\sigma^d)}}\right\}^{\beta},
\qquad
n^{-1/d}<\sigma\le L_n^{-1/2},
\]
with the notation defined in \cref{obj:def:frontier-rate}. Here the design-decay exponent enters through the logarithmic term as well as the first transition threshold. At a common sample size and error scale lying in the intermediate regime for the exponents being compared, a larger \(\kappa\) reduces \(n\sigma^d\) and increases the displayed uncertainty scale. The displayed scale depends on measurement precision, and once the direct regime is reached, the error-free weak-design order characterized in \cref{obj:prop:error-free-reduction} is attained. Further reductions within the direct regime share that minimax order.

Under the envelope calibration \(c_{\mathrm{lo}}\le(\kappa+1)2^\kappa\le c_{\mathrm{hi}}\), known compact support governs the eventual order when the measurement-error scale remains fixed and positive. Under \cref{obj:ass:latent-support}, every fixed \(\sigma\in(0,1/4]\) eventually lies in the compact-support regime as sample size increases. For this calibrated class, both optimal criteria then have order
\[
\left(\frac{\log\log n}{\log n}\right)^\beta,
\]
as established in \cref{obj:thm:uniform-frontier}. The design-decay exponent determines the direct and intermediate regimes, while smoothness determines the exponent in this eventual fixed-error order. Within the calibrated class, two fixed positive error scales therefore share this asymptotic order, with comparison constants allowed to depend on the error scale. The uniform characterization gives a more detailed account of improving precision with sample size: the compact-support expression depends on \(\sigma\) through \(\log(e+\sigma^2L_n)\), and sufficiently rapid improvement leads through the intermediate regime to the direct order.

Under the envelope calibration \(c_{\mathrm{lo}}\le(\kappa+1)2^\kappa\le c_{\mathrm{hi}}\), the matching estimation and interval-length orders concern uncertainty about the same scalar causal mean, \(\theta(P)\). Expected absolute error measures point-estimation accuracy, while honest expected length measures the width required for connected intervals with coverage of at least \(1-\alpha\) uniformly over the structural class; these criteria are defined in \cref{obj:def:minimax-risk,obj:def:honest-length}. For the calibrated class, their common rate connects improved measurement precision to both forms of uncertainty. The observable estimator and interval construction in \cref{obj:thm:observable-upper,obj:thm:finite-honesty} attain the corresponding upper bounds, and the observed-law comparison in \cref{obj:thm:observed-lower} supplies the converse under this calibration. Together, these results characterize the order of precision available for causal evaluation under the stated support, design, smoothness, and measurement conditions and the envelope calibration.
